\begin{figure}[H]
  \centering
  \begin{tikzpicture}
    \begin{axis}[paperaxis,
      xlabel={Année}, ylabel={Écart France moins comparables (part de l'actif)},
      xtick={2013,2014,2015,2016,2017,2018,2019,2020,2021,2022},
      xticklabel style={/pgf/number format/1000 sep=, rotate=45, anchor=east},
      legend pos=north west, legend cell align=left,
    ]
      \draw[zeroline, line width=0.6pt] (axis cs:2013,0) -- (axis cs:2022,0);
      % Deux positions de bilan de la France, toutes deux en bleu ; le style de trait distingue.
      \addplot[franceblue, line width=1.3pt, mark=*, mark options={scale=0.7}]
        table[col sep=comma, x=year, y=coef] {figures/cash_event.csv};
      \addlegendentry{Trésorerie sur actif}
      \addplot[franceblue, line width=1.3pt, densely dashed, mark=square*, mark options={scale=0.6}]
        table[col sep=comma, x=year, y=coef] {figures/fin_debt_event.csv};
      \addlegendentry{Dette financière sur actif}
      % Reforme de 2018 en vert ; entree dans la crise sanitaire en gris.
      \draw[baisse2018, dashed, line width=1pt]
        ({axis cs:2017.5,0}|-{rel axis cs:0,0}) -- ({axis cs:2017.5,0}|-{rel axis cs:0,1});
      \draw[controlgrey, dashed, line width=1pt]
        ({axis cs:2019.5,0}|-{rel axis cs:0,0}) -- ({axis cs:2019.5,0}|-{rel axis cs:0,1});
    \end{axis}
  \end{tikzpicture}
  \caption{Trésorerie et dette financière des sociétés françaises, écart aux pays comparables, autour de la crise sanitaire}
  \label{fig:covid-pge}
  \fignote{Coefficients $\tau_j$ de l'équation~\eqref{eq:event} pour la trésorerie sur actif et pour
  la dette financière, dette à long terme plus emprunts courants, sur actif, écart de la France aux
  cinq pays comparables l'année $j$, année de référence 2017 omise. Les deux séries sont en bleu, la
  couleur de la France dans tout le papier ; le trait plein est la trésorerie, le trait tireté la
  dette financière. La ligne verte verticale marque la réforme de 2018, la ligne grise l'entrée dans
  la crise sanitaire fin 2019. Les deux séries sont plates jusqu'en 2019 puis bondissent conjointement
  en 2020, 2021 et 2022, la signature des prêts garantis par l'État. Il s'agit d'une corrélation
  descriptive, non d'une preuve que la dette a financé la distribution.}
\end{figure}
